\documentclass[12pt]{article}
\newcommand{\bottomMargin}{2cm}

\input{structure.tex}

\begin{document}

\renewcommand{\tl}{$7.5$ сек.}
\renewcommand{\ml}{$256$ MB}

\problem{Задача AB12. Pointsort}
\addauthor{Автор: Емил Инджев}{2025}{03}{08}

\renewcommand{\header}{
    XLI НАЦИОНАЛНА ОЛИМмИАДА ПО ИНФОРМАТИКА\\
    Национален кръг\\
    Велико Търново, 7 - 10 март 2025 г.\\
    Група AB, 9-12 клас, Ден 1
}

\vspace{0.1in}

Алиса отново се оказва в Огледалния свят и отново се налага да сортира разни неща, които Хъмпти-Дъмпти е изтървал. Този път тя трябва да сортира $N$ $K$-измерни точки (да знае реда им по всяко от измеренията поотделно). За целта тя може да разглежда двойка точки и наведнъж да разбира коя е по-голяма и коя по-малка по всяко от измеренията. Тя иска да минимизира броя сравнения, които ще ползва. Отново обаче е без идеи и моли Вас за помощ.

По-формално, $N$-те точки са номерирани от $0$ до $N-1$. Всяка точка $i$ има $K$ на брой целочислени координати $X_{i,d}$ за $0 \leq d < K$. Също така знаем, че за всяко измерение $d$ координатите на точките по това измерение формират пермутация. Т.е. редицата $X_{0,d}, \dots, X_{N-1,d}$ е пермутация на числата на числата от $0$ до $N-1$. Гарантирано е, че тези пермутации са произволно избрани (с еднаква вероятност за всяка възможна пермутация) независимо една от друга.

Алиса може да сравни две точки $i$ и $j$ по всяко измерение наведнъж. Т.е. тя избира двете точки и като отговор получава списък от $K$ булеви резултата: $d$-тият резултат е единица, ако $X_{i,d} < X_{j,d}$, а иначе е нула. Целта на Алиса е чрез такива сравнения да разбере координатите на всички точки по всички измерения, т.е. стойностите на $X_{i,d}$ за всяко $i$ и $d$.

\subsection{Детайли по имплементацията}

Вие трябва да имплементирате функцията \verb|pointsort|:

\verb|std::vector<std::vector<int>> pointSort(int n, int k);|

Тя ще бъде извикана $T$ пъти с едни и същи $N$ и $K$ всеки път. Всяко извикване представлява независим събтест. Функцията трябва да върне стойностите на точките, индексирани първо по точки, а после по измерения.

Вашата програма може да вика функцията \verb|compare| на журито:

\verb|std::vector<bool> compare(int i, int j);|

Тя може да бъде извиквана произволна бройка пъти. Избраните $i$ и $j$ трябва да са различни един от друг валидни индекси на точки. Функцията връща резултатите на $K$-те сравнения по всяко от измеренията.

Вашата програма трябва да включва хедър файла \verb|pointsort.h|, не трябва да има \verb|main| функция и не трябва да чете от стандартния вход или да пише на стандартния изход.

\subsection{Ограничения}

\vspace{0.1em}

\begin{itemize}
\item $ N = 10^4 $
\item $ 2 \leq K \leq 5 $
\item $ 6 \leq T \leq 15 $
\end{itemize}

\subsection{Локално тестване}

Предоставени са Ви хедър файл и локален грейдър, които да компилирате с програмата си. Първо се въвеждат $T$, $N$ и $K$. След това се въвежда 1, ако искате локалния грейдър да генерира произволни събтестове, или 0, ако искате да въведете свои събтестове. След това се въвждат данните за всеки от $T$-те събтеста. Ако сме в режим 1, се въвежда само по едно число: сийдът за произволния генератор. Ако сме в режим 0, се въвеждат по $N$ реда от по $K$ числа: координатите на точките. Програмата ще изведе средния брой направени сравнения на събтест или съобщение за грешка при проблем.

\subsection{Оценяване}

Всеки от четирите теста на задачата се оценява поотделно. Резултатът Ви на даден тест (частта от точките за теста, която ще получите) се определя по следния начин:

Ако програмата Ви направи невалидно сравнение или върне грешен отговор, резултатът Ви е 0. Иначе, нека $Q$ е средният брой сравнения, които програмата Ви прави на събтест (т.е. средния брой сред $T$-те изпълнения). Също, нека $Q^*$ да е желания брой сравнения. Ако, $Q \leq Q^*$, резултатът Ви е 1. Иначе, дефинираме релативния брой допълнителни сравнения $E$ като:

$E = Q / Q^* - 1$

Тогава резултатът Ви $S$ е:

$S = \max(1 - 3.75 \times \min(E, 0.15) - \max(E - 0.15, 0), 0.1)$

Дадени са някои примерни стойности на $S$ като функция на $E$:

\begin{table}[H]
\begin{tblr}{|Q[c,m]|Q[c,m]|}
\hline
$E$ & $S$ \\
\hline
$0.00$ & $100.00\%$ \\
\hline
$0.05$ & $81.25\%$ \\
\hline
$0.10$ & $62.50\%$ \\
\hline
$0.15$ & $43.75\%$ \\
\hline
$0.20$ & $38.75\%$ \\
\hline
$0.30$ & $28.75\%$ \\
\hline
$0.40$ & $18.75\%$ \\
\hline
$0.50$ & $10.00\%$ \\
\hline
\end{tblr}
\end{table}

\subsection{Тестове}

\begin{table}[H]
\begin{tblr}{|Q[c,m]|Q[c,m]|Q[c,m]|Q[c,m]|Q[c,m]|}
\hline
\textbf{Тест} & \textbf{Точки} & $K$ & $T$ & $Q^*$ \\
\hline
$1$ & $25$ & $2$ & $15$ & $187500$ \\
\hline
$2$ & $35$ & $3$ & $10$ & $267000$ \\
\hline
$3$ & $25$ & $4$ & $8$ & $350800$ \\
\hline
$4$ & $15$ & $5$ & $6$ & $434500$ \\ 
\hline
\end{tblr}
\end{table}

\end{document}
